% Revisione Ready
% Grammatica Ready
% Citazioni Ready

\section{Analisi delle prestazioni computazionali}\label{sec:analisi_delle_presentazioni_computazionali}

In questa sezione vengono presentati e analizzati i riscontri empirici ottenuti dall'esecuzione della soluzione proposta (si veda il Capitolo~\ref{sec:soluzione_proposta}) mettendoli a confronto con delle versioni di confronto (si veda il Capitolo~\ref{sec:varianti_degli_algoritmi_analizzati}) sulle istanze di benchmark precedentemente caratterizzate (si veda il Capitolo~\ref{sec:istanze_di_benchmark_e_caratterizzazione_dei_dati}).
L'obiettivo fondamentale è quello di quantificare l'efficienza e l'efficacia dell'approccio basato sulla finestra temporale a scorrimento, evidenziandone la scalabilità al crescere della dimensionalità del problema.

\subsection{Metriche di valutazione e protocollo di tesi}

Per valutare qualitativamente il comportamento del sistema sono state analizzate le tre metriche precedentemente citate durante la definizione formale del problema (si veda il Capitolo~\ref{sec:definizione_formale}), ovvero:
\begin{itemize}
    \item \textbf{Tempo di esecuzione ($T_{exec}$)}: espresso in secondi, misura il tempo impiegato da parte dell'elaboratore per arrivare al termine dell'esecuzione dell'algoritmo analizzato.
    \item \textbf{Dimensione media dell'Hitting Set ($\overline{|X|}_{\tau}$)}: espresso attraverso il numero di nodi, misura la cardinalità media temporale dell'insieme di copertura calcolato durante l'intero periodo di esecuzione.
    \item \textbf{Numero nodi distinti scelti ($R$)}: espresso attraverso il numero di nodi, rappresenta il volume complessivo di vertici unici accumulati nell'Hitting Set globale al termine del processo (si veda il Capitolo~\ref{sec:definizione_formale}).
\end{itemize}
L'analisi congiunta delle tre metriche permette di mappare accuratamente il trade-off dell'algoritmo proposto tra l'efficienza computazionale del codice e la qualità delle soluzioni generate, offrendo una panoramica completa sulla scalabilità dell'approccio proposto.

\subsection{Aspettative sperimentali e risultati}

Prima di analizzare direttamente i risultati sperimentali, è importante definire un comportamento atteso dell'algoritmo in tutte le sue versioni in modo da poter successivamente confrontare le aspettative con i risultati effettivi.
Poiché l'efficienza della soluzione si basa parzialmente anche sul rapporto tra la dimensione della finestra temporale considerata $(\epsilon)$ e la dimensione dello scorrimento computazionale della stessa $(\delta)$, abbiamo deciso, per ogni dataset analizzato, di testare tre rapporti differenti tra $\epsilon$ e $\delta$ in modo da poter analizzare come cambia il risultato al variare delle loro dimensioni.

\subsubsection*{Contacts high school}

La caratteristica particolare di questo dataset è la distribuzione di arrivo degli iperarchi marcatamente intermittente (si veda Figura~\ref{fig:distribuzione_dataset} a).
Su tale base si può affermare che l'efficacia dell'algoritmo proposto e delle rispettive estensioni può essere garantita solo in base al rapporto che la dimensione della finestra temporale e la dimensione del relativo scorrimento hanno con i picchi e le valli di attività della distribuzione mostrata.
La soluzione proposta riesce dunque a garantire prestazioni ottimali solo a condizione che l'ampiezza della finestra temporale sia sufficiente a intercettare più picchi concorrenti e che il passo di scorrimento sia almeno pari al loro tempo di inter-arrivo.
In pratica l'unico modo affinché l'algoritmo riesca a funzionare con questa distribuzione degli archi in ingresso è selezionando delle dimensioni di $\epsilon$ e $\delta$ in grado di mascherare la presenza dei burst temporali.

\begin{figure}[htbp]
    \centering
    \begin{minipage}{0.32\textwidth}
        \centering
        \includegraphics[width=\textwidth, trim={0cm 0cm 5cm 0cm}, clip]{Immagini/time_distribution_graphs/contacts_high_school.pdf}
        \vspace{-9mm}
        \caption*{Tempo di esecuzione}
    \end{minipage}
    \hfill
    \begin{minipage}{0.32\textwidth}
        \centering
        \includegraphics[width=\textwidth, trim={0cm 0cm 5cm 0cm}, clip]{Immagini/size_distribution_graphs/contacts_high_school.pdf}
        \vspace{-9mm}
        \caption*{Dimensione media Hitting Set}
    \end{minipage}
    \hfill
    \begin{minipage}{0.32\textwidth}
        \centering
        \includegraphics[width=\textwidth, trim={0cm 0cm 5cm 0cm}, clip]{Immagini/reuse_distribution_graphs/contacts_high_school.pdf}
        \vspace{-9mm}
        \caption*{Numero nodi distinti scelti}
    \end{minipage}

    \vspace{2mm}
    % Legenda manuale
    
\begin{tikzpicture}
        \draw[blue, thick] (0,0) -- (0.5,0); \node[right] at (0.5,0) {Iterative Greedy};
        \draw[green, thick] (5,0) -- (5.5,0); \node[right] at (5.5,0) {Sliding different Greedy};
        \draw[orange, thick] (10,0) -- (10.5,0); \node[right] at (10.5,0) {Sliding New Nodes};
        \draw[red, thick] (15,0) -- (15.5,0); \node[right] at (15.5,0) {Sliding Reuse};
    \end{tikzpicture}
    
\end{figure}

\subsubsection*{Coauth history}

Come nel grafico precedente, la caratteristica particolare del dataset risiede nella sua distribuzione di arrivo degli iperarchi (si veda Figura~\ref{fig:distribuzione_dataset} b).
La presenza di un picco di nodi verso la fine del periodo temporale studiato non permette di sfruttare la conoscenza acquisita dalle finestre precedenti, limitando le situazioni in cui l'algoritmo riesce a fornire prestazioni migliori.

\begin{figure}[htbp]
    \centering
    \begin{minipage}{0.32\textwidth}
        \centering
        \includegraphics[width=\textwidth, trim={0cm 0cm 5cm 0cm}, clip]{Immagini/time_distribution_graphs/coauth_history.pdf}
        \vspace{-9mm}
        \caption*{Tempo di esecuzione}
    \end{minipage}
    \hfill
    \begin{minipage}{0.32\textwidth}
        \centering
        \includegraphics[width=\textwidth, trim={0cm 0cm 5cm 0cm}, clip]{Immagini/size_distribution_graphs/coauth_history.pdf}
        \vspace{-9mm}
        \caption*{Dimensione media Hitting Set}
    \end{minipage}
    \hfill
    \begin{minipage}{0.32\textwidth}
        \centering
        \includegraphics[width=\textwidth, trim={0cm 0cm 5cm 0cm}, clip]{Immagini/reuse_distribution_graphs/coauth_history.pdf}
        \vspace{-9mm}
        \caption*{Numero nodi distinti scelti}
    \end{minipage}

    \vspace{2mm}
    % Legenda manuale
    
\begin{tikzpicture}
        \draw[blue, thick] (0,0) -- (0.5,0); \node[right] at (0.5,0) {Iterative Greedy};
        \draw[green, thick] (5,0) -- (5.5,0); \node[right] at (5.5,0) {Sliding different Greedy};
        \draw[orange, thick] (10,0) -- (10.5,0); \node[right] at (10.5,0) {Sliding New Nodes};
        \draw[red, thick] (15,0) -- (15.5,0); \node[right] at (15.5,0) {Sliding Reuse};
    \end{tikzpicture}
    
\end{figure}

\subsubsection*{Email Enron}

A differenza dei grafici precedenti, la distribuzione di arrivo degli iperarchi del dataset non sembra poter creare particolari problemi, in quanto la fase di picco è situata a metà del periodo temporale studiato (si veda Figura~\ref{fig:distribuzione_dataset} c).
Tale dataset però risulta essere particolarmente denso (si veda Tabella~\ref{tab:caratteristiche_dataset}), fattore che potrebbe portare a una minore differenza nel numero di nodi distinti scelti $(R)$ e a un buon grado di approssimazione dell'Hitting Set sia da parte degli algoritmi basati su finestra temporale a scorrimento, sia da parte di quelli greedy iterativi.

\begin{figure}[htbp]
    \centering
    \begin{minipage}{0.32\textwidth}
        \centering
        \includegraphics[width=\textwidth, trim={0cm 0cm 5cm 0cm}, clip]{Immagini/time_distribution_graphs/email_Enron.pdf}
        \vspace{-9mm}
        \caption*{Tempo di esecuzione}
    \end{minipage}
    \hfill
    \begin{minipage}{0.32\textwidth}
        \centering
        \includegraphics[width=\textwidth, trim={0cm 0cm 5cm 0cm}, clip]{Immagini/size_distribution_graphs/email_Enron.pdf}
        \vspace{-9mm}
        \caption*{Dimensione media Hitting Set}
    \end{minipage}
    \hfill
    \begin{minipage}{0.32\textwidth}
        \centering
        \includegraphics[width=\textwidth, trim={0cm 0cm 5cm 0cm}, clip]{Immagini/reuse_distribution_graphs/email_Enron.pdf}
        \vspace{-9mm}
        \caption*{Numero nodi distinti scelti}
    \end{minipage}

    \vspace{2mm}
    % Legenda manuale
    
\begin{tikzpicture}
        \draw[blue, thick] (0,0) -- (0.5,0); \node[right] at (0.5,0) {Iterative Greedy};
        \draw[green, thick] (5,0) -- (5.5,0); \node[right] at (5.5,0) {Sliding different Greedy};
        \draw[orange, thick] (10,0) -- (10.5,0); \node[right] at (10.5,0) {Sliding New Nodes};
        \draw[red, thick] (15,0) -- (15.5,0); \node[right] at (15.5,0) {Sliding Reuse};
    \end{tikzpicture}
    
\end{figure}

\subsubsection*{Crypto bitcoin}

L'ultimo dataset risulta essere quello più favorevole dal punto di vista dell'algoritmo da me sviluppato, in quanto presenta una distribuzione di arrivo quasi costante e di una struttura non particolarmente densa.
In presenza di tali condizioni, ci aspettiamo che il nostro algoritmo riesca a sfruttare al meglio il funzionamento della finestra temporale, mentre, a causa della natura sparsa del dataset, ci aspettiamo che gli algoritmi greedy iterativi non riescano a ottenere risultati altrettanto buoni.

\begin{figure}[htbp]
    \centering
    \begin{minipage}{0.32\textwidth}
        \centering
        \includegraphics[width=\textwidth, trim={0cm 0cm 5cm 0cm}, clip]{Immagini/time_distribution_graphs/crypto_bitcoin_201601.pdf}
        \vspace{-9mm}
        \caption*{Tempo di esecuzione}
    \end{minipage}
    \hfill
    \begin{minipage}{0.32\textwidth}
        \centering
        \includegraphics[width=\textwidth, trim={0cm 0cm 5cm 0cm}, clip]{Immagini/size_distribution_graphs/crypto_bitcoin_201601.pdf}
        \vspace{-9mm}
        \caption*{Dimensione media Hitting Set}
    \end{minipage}
    \hfill
    \begin{minipage}{0.32\textwidth}
        \centering
        \includegraphics[width=\textwidth, trim={0cm 0cm 5cm 0cm}, clip]{Immagini/reuse_distribution_graphs/crypto_bitcoin_201601.pdf}
        \vspace{-9mm}
        \caption*{Numero nodi distinti scelti}
    \end{minipage}

    \vspace{2mm}
    % Legenda manuale
    
\begin{tikzpicture}
        \draw[blue, thick] (0,0) -- (0.5,0); \node[right] at (0.5,0) {Iterative Greedy};
        \draw[green, thick] (5,0) -- (5.5,0); \node[right] at (5.5,0) {Sliding different Greedy};
        \draw[orange, thick] (10,0) -- (10.5,0); \node[right] at (10.5,0) {Sliding New Nodes};
        \draw[red, thick] (15,0) -- (15.5,0); \node[right] at (15.5,0) {Sliding Reuse};
    \end{tikzpicture}
    
\end{figure}